\section{External Diagnostics and Reproducible Workflow}
\label{app:external}

External diagnostics evaluate whether model outputs are consistent with specified
evidence outside the fitted likelihood. In this appendix, ``consistent'' should be
defined against explicit data sources or biological constraints, rather than by informal
judgement. Generic code for reading, plotting, and comparing indicators should live in
packages; case-study scripts should map model outputs to external evidence and write
results into the content files below.

\subsection{Required R Packages}

\begin{table}

\caption{\label{tab:external-libraries}Required R packages for external diagnostics.}
\centering
\begin{tabular}[t]{llll}
\toprule
package & role & tier & installed\\
\midrule
knitr & Knit template; format LaTeX tables & core & 1.51\\
r4ss & Read SS3 output for maturity, selectivity, and index fit & run\_code & 1.53.0\\
sstruth & diagSS, spSsbFromSs3, ssPeDiagnose; external CPUE and SP plots & run\_code & 0.0.0.9020\\
FLRebuild & curveSS(); production and reference points from SS3 & run\_code & 0.1.0\\
ss3om & readTruth / readFLSss3(); FLR stock objects from SS3 & run\_code & 0.5.3.9004\\
\addlinespace
FLCandy & matSel(); maturity vs selectivity ogives & run\_code & 1.7.1\\
diags & Extract CPUE/index series from SS3 diagnostics (.diagUss) & run\_code & 2.2.1\\
mpb & Lowess smoothing and scaling for CPUE time-series plots & run\_code & 3.5.0\\
plyr & Reshape index data (ddply) & run\_code & 1.8.9\\
reshape2 & Wide format for correlation plots (cast) & run\_code & 1.4.5\\
\addlinespace
ggplot2 & CPUE time series, survey vs biomass, sel/mat, LBI & run\_code & 4.0.3\\
GGally & Pairwise scatter plots of index correlations (planned) & planned & 2.4.0\\
ggcorrplot & Correlation matrix with hierarchical clustering (planned) & planned & 0.1.4.1\\
FLCore & FLR indicators and survey standardisation objects (planned) & planned & 2.6.29.9001\\
FLife & Cross-correlations between indices (planned) & planned & 3.4.0\\
\addlinespace
mydas & Life-history indicators (LBI); my\_smooth for ggpairs (planned) & planned & 1.2.2\\
\bottomrule
\end{tabular}
\end{table}



\subsection{Evidence Outside the Likelihood}

List the external evidence used in the case study:

\begin{itemize}
  \item life-history information, including growth, maturity, natural mortality, and longevity;
  \item fishery-independent survey trends, spatial indicators, and availability information;
  \item fishery behaviour, selectivity evidence, discard or bycatch information;
  \item environmental or ecosystem drivers already used in index standardisation;
  \item demographic indicators such as mature fraction, older-fish indicators, or age structure.
\end{itemize}

\begin{table}

\caption{\label{tab:external-evidence}External evidence register for the shortfin mako SS3 case study.}
\centering
\begin{tabular}[t]{lllll}
\toprule
diagnostic & evidence\_source & model\_quantity & criterion & status\\
\midrule
Growth and maturity & Life-history inputs to SS3 & Maturity-at-age vs selectivity-at-age & Fishing should not heavily precede 50 percent maturity on SSB scale & Check\\
Selectivity at age & Fleet selectivity ogives from SS3 & Selectivity ogive & Selectivity peak below maturity implies lagged SSB response & Mismatch flagged in case study\\
Survey perception & Standardised CPUE / survey indices (external to likelihood if pre-processed) & Index fit and scaling & Assessment trend coherent with survey perception & Check\\
Age structure & Length and age composition data & Age composition residuals & Older-fish / age-structure signals consistent with depletion narrative & Check\\
Spatial availability & Spatial models / availability (if used outside SS3) & Recruitment / availability deviations & No double-counting of spatial signal inside SS3 & Check\\
\addlinespace
Fishery behaviour / non-retention & Discard and non-retention regulations & Catch-at-age and discard mortality & Removals dominated by juveniles while state is mature female SSB & Mismatch flagged in case study\\
\bottomrule
\end{tabular}
\end{table}



\subsection{Standard Checks}

For SS3 and other age-structured assessments, useful external checks include maturity
versus selectivity, survey trend versus assessment biomass, spatial occupancy or aggregation
indicators, and comparison of recruitment or productivity regimes with external evidence.

% External comparison tables.

\begin{table}

\caption{\label{tab:external-sspe}Surplus-production vs SS3 process-error diagnostics (sstruth::ssPeDiagnose).}
\centering
\begin{tabular}[t]{lrrrrrrrl}
\toprule
scenario & curve\_cor & curve\_rmse & curve\_rel\_rmse & pe\_cor\_ssb & pe\_cor\_biomass & pe\_rmse\_ssb & pe\_rmse\_biomass & likely\_misspecified\\
\midrule
Scenario 1 & 0.9933 & 374.4 & 0.1159 & -0.2136 & 0.2427 & 7.196 & 0.006521 & TRUE\\
\bottomrule
\end{tabular}
\end{table}



\subsection{Diagnostic Output}

Tables, figures, and code-generated summaries can be added to the file below as they
become available.

% External diagnostic figures.

\begin{figure}[H]
  \centering
  \includegraphics[width=0.72\linewidth]{appendices/figs/external/cpue-ts.png}
  \caption{Standardised CPUE time series by fleet (sstruth::diagSS; lowess trend when mpb available).}
  \label{fig:external-cpue-ts}
\end{figure}

\begin{figure}[H]
  \centering
  \includegraphics[width=0.88\linewidth]{appendices/figs/external/sp-ssb.png}
  \caption{Empirical surplus production vs spawning stock biomass (sstruth::spSsbFromSs3).}
  \label{fig:external-sp-ssb}
\end{figure}

\begin{figure}[H]
  \centering
  \includegraphics[width=0.88\linewidth]{appendices/figs/external/selmat.png}
  \caption{Comparison of selectivity and maturity-at-age in the shortfin mako SS3 scenarios (from readFLSss3).}
  \label{fig:external-selmat}
\end{figure}



\subsection{Interpretation}

State which model outputs are consistent with external evidence, which are not, and which
external conflicts should be tested in the robustness appendix.

% Case-specific external interpretation.
% State which model outputs are consistent with life history and fishery evidence,
% which conflicts should be tested in the robustness appendix,
% and whether the biomass currency matches the age classes removed by fishing.


